\documentclass[en]{snu-ece-bsc-thesis}

% Add your packages here, e.g.,
% \usepackage{tikz}
\usepackage{siunitx}

% For lorem ipsum; remove these lines when writing your thesis
\usepackage{lipsum}
\usepackage{jiwonlipsum}

% hyperref *must* be the last package to be loaded!
\usepackage[pdfusetitle]{hyperref}

\addbibresource{bib.bib}

\title{SNU ECE BSc Thesis Template}
\author{이재호}
\advisor{홍길동}
\date{2024년 2월}
\approvaldate{2023년 12월 15일}

\koreankeywords{서울대학교, 공과대학, 학사학위논문, 템플릿}
\englishkeywords{Seoul National University, College of Engineering, Bachelor's Thesis, Template}


\begin{document}
\maketitle

\pagenumbering{roman}
\begin{abstract}
  \lipsum[1]
\end{abstract}

\tableofcontents
\listoftables
\listoffigures

\chapter{Introduction}\label{chap:introduction}
\pagenumbering{arabic}
This template is structured as follows.
Section~\ref{sec:figure} of chapter~\ref{chap:body} shows an example of a figure.
Section~\ref{sec:table} of chapter~\ref{chap:body} shows an example of a table.
Chapter~\ref{chap:conclusion} concludes this template.
감격스럽다
% \lipsum[2-3]


\chapter{Body}\label{chap:body}
Entropy of information is the expected value of information contained in each message, and is given by equation~\eqref{eq:entropy}~\cite{6773024}.
\begin{equation}\label{eq:entropy}
  H(X) = -\sum_{i=1}^n {\mathrm{P}(x_i) \log_b \mathrm{P}(x_i)}
\end{equation}

\lipsum[4-6]


\section{Figure}\label{sec:figure}
Example of a figure is shown in figure~\ref{fig:example}.
Figure~\ref{fig:snu} is the logo of Seoul National University and figure~\ref{fig:eng} is the logo of College of Engineering.

\begin{figure}[htp]
  \centering
  \begin{subfigure}[b]{0.5\textwidth}
    \centering
    \includegraphics[width=0.5\textwidth]{logo1.pdf}
    \caption{The logo of Seoul National University}\label{fig:snu}
  \end{subfigure}%
  \begin{subfigure}[b]{0.5\textwidth}
    \centering
    \includegraphics[width=0.9\textwidth]{logo2.pdf}
    \caption{The logo of College of Engineering}\label{fig:eng}
  \end{subfigure}
  \caption[Figure example (ToC)]{An example of a figure.}\label{fig:example}
\end{figure}

\lipsum[7-8]


\section{Table}\label{sec:table}
Example of a table is shown in table~\ref{tab:example}.\footnote{\lipsum[12]}

\begin{table}[htp]
  \centering
  \caption[Table example (ToC)]{An example of a table.}\label{tab:example}
  \begin{tblr}{cc}
    \toprule
    Constant & Value \\\midrule
    $c$ & \SI{299792458}{\meter\per\second} \\
    $h$ & \SI{6.62607015e-34}{\joule\per\hertz} \\\bottomrule
  \end{tblr}
\end{table}

\lipsum[9-10]


\chapter{Conclusion}\label{chap:conclusion}
\lipsum[11]

\printbibliography

\begin{abstract}[ko]
  \jiwon[1]
\end{abstract}
\end{document}
